\section{Computability and the Refining Universe}

A discrete, local base is a computable base. This matters, and it is easy to overstate, so here is the careful
version. Computability is the property that a process could in principle be run, step by step, by an idealized
computer. A supertask is the completion of infinitely many steps in a finite time, which no machine can do. With
these two terms fixed, the consequences for the substrate are sharp and limited.

\subsection{What computability does and does not forbid}

Computability forbids the continuum, the uncomputable real numbers, an actual smooth infinity of points between
any two points. It forbids supertasks, the finishing of infinitely many steps to get one answer. It does not
forbid an infinite universe. An infinite mesh is fine, as long as each finite region is finite and each step is
local and finite. Size is not the issue. Continuity and supertasks are.

This corrects an earlier overclaim. Computability does not push the theory toward finite and growing. It pushes
the theory toward discrete and local. Both the finite universe and the infinite universe pass, because both are
discrete. The fork between a finite beginning and an eternal past is left untouched by the computability argument,
which constrains the texture of the base but not the length of its history.

\subsection{The load-bearing commitment}

\begin{principle}[Discrete, not continuum]
The base is discrete in every respect. There is a mesh, $24$ directions, a ternary tone, a discrete beat, and a
finite symmetry group. Continuity is never part of the base. The Lie groups, the Dirac field, smooth Lorentz
invariance, and the $1/r$ law are always emergent, a coarse-grained limit of the discrete weave seen from far
away. The universe is a computation, and the smoothness we observe is the look of a fine-enough discrete fabric at
a distance.
\end{principle}

This commitment is what the rest of the framework rests on. It is also what separates this program from the many
proposals that take a continuous manifold as primitive and then quantize it. Here the order of explanation runs
the other way. The discrete structure is given, and the continuum is what the structure looks like when you stop
resolving cells.

\subsection{The refining picture and its real limit}

It is tempting to say the universe refines toward the continuum as it grows. The precise version is narrower. The
large-scale, infrared description looks continuous because the discrete steps are tiny relative to us, not because
the base ever becomes continuous. The base stays discrete forever. The continuum is a viewpoint, not a
destination. The universe does not approach smoothness over time. It is discrete at every scale of its own
construction, and smoothness is only ever a feature of the observer's coarse view.


A computable cosmos means the theory is in principle fully simulable. There is no magic and no uncomputable
ingredient. It also frees the infinite-universe story, because infinity was never the problem. The continuum was.
An infinite past and an infinite mesh both survive the computability test, while the smooth continuum and the
supertask do not. In the carrying image, the fire is made of countably many sparks, however far it spreads.
